\section{Methodology}

To establish the relationships among the selected languages, we evaluated them
across three dimensions: orthographic, phonetic, and geospatial.

\subsection{Linguistic Distance Metrics}

We extracted unique trigrams from the standardized biblical corpus of each
language to compute two distinct similarity matrices.

First, we measured orthographic similarity using the Jaccard similarity
coefficient, which quantifies the direct overlap of written forms. It is
defined as the intersection of two trigram sets divided by their union:

\begin{equation}
    J(A, B) = \frac{|A \cap B|}{|A \cup B|}
\end{equation}

where \(A\) and \(B\) represent the unique orthographic trigrams for two
languages.

Second, we measured phonetic similarity to capture acoustic properties
independent of spelling conventions. We converted all corpus trigrams into
phonetic representations using the Double Metaphone algorithm, which accounts
for non-English phonological patterns and variant pronunciations. We then
computed the pairwise cosine similarity of these phonetic frequency
distributions:

\begin{equation}
    \text{cos}(\theta) = 
        \frac{\mathbf{A} \cdot \mathbf{B}}{\|\mathbf{A}\| \|\mathbf{B}\|} = 
        \frac{\sum_{i=1}^{n} A_i B_i}{\sqrt{\sum_{i=1}^{n} A_i^2} \sqrt{\sum_{i=1}^{n} B_i^2}}
\end{equation}

where \(\mathbf{A}\) and \(\mathbf{B}\) are the phonetic frequency vectors.
Cosine similarity accounts for the relative distribution of features rather
than absolute frequencies, making it robust to corpus size variations.

\subsection{Geospatial Mapping}

To model how physical terrain influences these linguistic similarities, we
aligned the linguistic data with geospatial boundaries. We sourced language
distribution boundaries from the Komisyon ng Wikang Filipino (KWF)
atlases~\cite{Marrion2025MapaKultural, KWF2025MapaPilipinas} and administrative
data from the UN Office for the Coordination of Humanitarian
Affairs~\cite{OCHA_Philippines_2018_PHL_Admin_Boundaries}. Topographic data
(elevation and hydrology) were obtained from the National Mapping and Resource
Information Authority (NAMRIA)~\cite{NAMRIA_Downloads}. We rendered the
combined digital map using Matplotlib and GeoPandas in
Python~\cite{Hunter:2007, kelsey_jordahl_2020_3946761} to visually and
quantitatively assess the convergence of linguistic boundaries and terrain
discontinuities.